% GENERATED FILE. Edit canonical lessons, then run npm run generate:books.
% canonical-source-hash: fnv1a64:cfb37e9c7702fb3e
% canonical-lessons: LA-C365-quandoquidem, LA-C365-quippe, LA-C365-etenim, LA-C365-nempe, LA-C365-scilicet

\chapter{For, Of Course, Namely}
\label{ch:la-for-of-course-namely}
\hlchaptermodality{365}

\begin{chapteropening}
\textbf{By the end of this chapter:} \emph{I can back up a claim with since, for, of course and namely, and say what mood each word takes.}

Complete the last lesson of chapter 365: I can back up a claim with since, for, of course and namely, and say what mood each word takes.
\end{chapteropening}

\section[quandōquidem]{quandōquidem — since, seeing that}

\label{lesson:LA-C365-quandoquidem}



\begin{quote}
Before the new one: say the Latin for to indicate, then the Latin for a suspicion.
\end{quote}

\subsection*{You'll want to know: quandōquidem}
\textbf{quandōquidem} — \textquotedblleft{}since, seeing that\textquotedblright{}.

\begin{quote}
Quandōquidem tū ades, incipiāmus.
\end{quote}

\begin{quote}
\emph{Since you are here, let us begin.}
\end{quote}

\emph{Quandōquidem} takes the indicative: the reason is a plain fact both sides accept. It is built from \emph{quandō} and \emph{quidem}, indeed.

\begin{cousinweb}
\emph{quandōquidem} is built from \emph{quandō} \textquotedblleft{}when, since\textquotedblright{} and the particle \emph{quidem} \textquotedblleft{}indeed\textquotedblright{} (Lewis \& Short): since, as we both know.
\end{cousinweb}

\begin{grammarlens}[title={What you've built}]
One more word for telling, explaining and describing at B1.
\end{grammarlens}

\subsection*{Your turn}
\begin{itemize}
\raggedright
  \item \emph{Say it:} \emph{quandōquidem}
  \item \emph{Say it:} \emph{quandōquidem}, once more
  \item \emph{Say it:} a sentence of your own with \emph{quandōquidem}
  \item \emph{Recall:} say the Latin for to indicate, then the Latin for a suspicion, then say \emph{quandōquidem} again
\end{itemize}

\begin{tcolorbox}[breakable,colback=teal!4,colframe=teal!35!black,title={Before you move on}]
What does quandōquidem mean? (\textquotedblleft{}Since, seeing that\textquotedblright{}.) Say it once more.
\end{tcolorbox}
\section[quippe]{quippe — of course, naturally; for}

\label{lesson:LA-C365-quippe}



\begin{quote}
Before the new one: say the Latin for a suspicion, then the Latin for since.
\end{quote}

\subsection*{You'll want to know: quippe}
\textbf{quippe} — \textquotedblleft{}of course, naturally; for\textquotedblright{}.

\begin{quote}
Nōn vēnit, quippe quī aeger esset.
\end{quote}

\begin{quote}
\emph{He did not come, since of course he was ill.}
\end{quote}

\emph{Quippe} gives a reason as something obvious, sometimes with irony. \emph{Quippe quī} + subjunctive, seeing that he, is the classical use; Plautus and Livy also put the indicative after it.

\begin{cousinweb}
\emph{quippe} is built from the old adverb \emph{quī} \textquotedblleft{}how\textquotedblright{} and the particle \emph{-pe} (Lewis \& Short), the same ending that is in \emph{nempe}.
\end{cousinweb}

\begin{grammarlens}[title={What you've built}]
One more word for telling, explaining and describing at B1.
\end{grammarlens}

\subsection*{Your turn}
\begin{itemize}
\raggedright
  \item \emph{Say it:} \emph{quippe}
  \item \emph{Say it:} \emph{quippe}, once more
  \item \emph{Say it:} a sentence of your own with \emph{quippe}
  \item \emph{Recall:} say the Latin for a suspicion, then the Latin for since, then say \emph{quippe} again
\end{itemize}

\begin{tcolorbox}[breakable,colback=teal!4,colframe=teal!35!black,title={Before you move on}]
What does quippe mean? (\textquotedblleft{}Of course, naturally; for\textquotedblright{}.) Say it once more.
\end{tcolorbox}
\section[etenim]{etenim — for indeed, and in fact}

\label{lesson:LA-C365-etenim}



\begin{quote}
Before the new one: say the Latin for since, then the Latin for of course.
\end{quote}

\subsection*{You'll want to know: etenim}
\textbf{etenim} — \textquotedblleft{}for indeed, and in fact\textquotedblright{}.

\begin{quote}
Domī mānsī; etenim pluēbat.
\end{quote}

\begin{quote}
\emph{I stayed at home: for it was raining.}
\end{quote}

\emph{Etenim} (= \emph{et} + \emph{enim}) opens a sentence that gives the reason for the last one. Like \emph{enim} and \emph{nam}, it takes the indicative.

\begin{cousinweb}
\emph{etenim} is built from \emph{et} \textquotedblleft{}and\textquotedblright{} and \emph{enim} \textquotedblleft{}for, you see\textquotedblright{} (Lewis \& Short): and the reason is.
\end{cousinweb}

\begin{grammarlens}[title={What you've built}]
One more word for telling, explaining and describing at B1.
\end{grammarlens}

\subsection*{Your turn}
\begin{itemize}
\raggedright
  \item \emph{Say it:} \emph{etenim}
  \item \emph{Say it:} \emph{etenim}, once more
  \item \emph{Say it:} a sentence of your own with \emph{etenim}
  \item \emph{Recall:} say the Latin for since, then the Latin for of course, then say \emph{etenim} again
\end{itemize}

\begin{tcolorbox}[breakable,colback=teal!4,colframe=teal!35!black,title={Before you move on}]
What does etenim mean? (\textquotedblleft{}For indeed, and in fact\textquotedblright{}.) Say it once more.
\end{tcolorbox}
\section[nempe]{nempe — surely, of course; namely}

\label{lesson:LA-C365-nempe}



\begin{quote}
Before the new one: say the Latin for of course, then the Latin for for indeed.
\end{quote}

\subsection*{You'll want to know: nempe}
\textbf{nempe} — \textquotedblleft{}surely, of course; namely\textquotedblright{}.

\begin{quote}
Quis hoc fēcit? Nempe tū.
\end{quote}

\begin{quote}
\emph{Who did this? You, of course.}
\end{quote}

\emph{Nempe} appeals to what the listener must grant: surely, of course. It takes the indicative.

\begin{cousinweb}
\emph{nempe} is built from \emph{nam} \textquotedblleft{}for\textquotedblright{} and the particle \emph{-pe} (EtymDB), the same ending that is in \emph{quippe}.
\end{cousinweb}

\begin{grammarlens}[title={What you've built}]
One more word for telling, explaining and describing at B1.
\end{grammarlens}

\subsection*{Your turn}
\begin{itemize}
\raggedright
  \item \emph{Say it:} \emph{nempe}
  \item \emph{Say it:} \emph{nempe}, once more
  \item \emph{Say it:} a sentence of your own with \emph{nempe}
  \item \emph{Recall:} say the Latin for of course, then the Latin for for indeed, then say \emph{nempe} again
\end{itemize}

\begin{tcolorbox}[breakable,colback=teal!4,colframe=teal!35!black,title={Before you move on}]
What does nempe mean? (\textquotedblleft{}Surely, of course; namely\textquotedblright{}.) Say it once more.
\end{tcolorbox}
\section[scīlicet]{scīlicet — to be sure, of course}

\label{lesson:LA-C365-scilicet}



\begin{quote}
Before the new one: say the Latin for for indeed, then the Latin for surely.
\end{quote}

\subsection*{You'll want to know: scīlicet}
\textbf{scīlicet} — \textquotedblleft{}to be sure, of course\textquotedblright{}.

\begin{quote}
Tū scīlicet omnia scīs!
\end{quote}

\begin{quote}
\emph{You know everything, to be sure!}
\end{quote}

\emph{Scīlicet} = \emph{scīre licet}, one may know. Often ironic: to be sure! English \emph{scilicet}, short \emph{sc.}, means namely.

\begin{cousinweb}
\emph{scīlicet} is contracted from \emph{scīre licet} \textquotedblleft{}one may know\textquotedblright{} (Lewis \& Short): it is plain to know. English \emph{scilicet} \textquotedblleft{}namely\textquotedblright{} was borrowed from it.
\end{cousinweb}

\begin{grammarlens}[title={What you've built}]
One more word for telling, explaining and describing at B1.
\end{grammarlens}

\subsection*{Your turn}
\begin{itemize}
\raggedright
  \item \emph{Say it:} \emph{scīlicet}
  \item \emph{Say it:} \emph{scīlicet}, once more
  \item \emph{Say it:} a sentence of your own with \emph{scīlicet}
  \item \emph{Recall:} say the Latin for for indeed, then the Latin for surely, then say \emph{scīlicet} again
\end{itemize}

\begin{tcolorbox}[breakable,colback=teal!4,colframe=teal!35!black,title={Before you move on}]
What does scīlicet mean? (\textquotedblleft{}To be sure, of course\textquotedblright{}.) Say it once more.
\end{tcolorbox}
